\section{两场讲座，一个研究方法：先看数据，再研究变化}

这节课由两场风格不同、逻辑互补的讲座组成。Jason Wei 从训练数据和行为曲线出发，追问一个简单的 next-token objective 为什么会产生如此多能力；Hyung Won Chung 则从 Transformer 的历史变化出发，追问哪些结构是当前算力条件下的捷径，哪些又会在规模扩大后变成瓶颈。把两场讲座放在一起，得到的不是两组零散观点，而是一套共同方法：\textbf{先接触真实证据，再把观察分解成可检验机制，最后沿着规模或时间轴研究趋势。}

\subsection{Jason 的问题：为什么语言模型会工作得这么好}

本节先建立第一场讲座的研究姿态。标题页没有承诺一个完整理论，而是明确说要给出若干 intuition。这里的 intuition 不是随意猜测，而是从数据、损失曲线和具体 prompt 中压缩出的工作假设；它们必须能指导下一次实验，而不能只在事后解释成功。

\lecturefigure{jason-slide-001.jpg}{Jason Wei 的讲座主题：Intuitions on Language Models}{Jason Wei official deck, p.~1}

\begin{importantbox}{第一场讲座的四个问题}
\begin{enumerate}
  \item next-word prediction 为什么不只是一个任务，而像海量任务的混合？
  \item 规模扩大为什么能稳定降低总体 loss？
  \item 总体 loss 平滑时，单个任务为什么会呈现突变？
  \item 为什么某些精心挑选的任务会出现 inverse 或 U-shaped scaling？
\end{enumerate}
\end{importantbox}

\lecturefigure{jason-slide-002.jpg}{基本问题与手工检查训练数据的方法}{Jason Wei official deck, p.~2}

读这页时要先看右下角的方法，而不是只看标题。``Why do LLMs work?'' 太大，无法直接实验；``manually inspect data'' 则把问题变成可以执行的观察：随机抽样文本、手工做 continuation、标出一个 token 需要调用的知识与规则，再问模型可能从哪些上下文学到这些能力。它把抽象模型研究重新连接到训练分布。

\teachervoice{Jason 在 00:00:55--00:01:27 把“手工看数据”称为一个极其有帮助的工具。这个提醒反对只看 benchmark 聚合分数：如果研究者从未读过训练样本，就很容易把数据中显而易见的监督信号误写成神秘涌现。}

\lecturefigure{jason-slide-003.jpg}{手工看数据是在训练研究者自己的“生物神经网络”}{Jason Wei official deck, p.~3}

这句话的教学含义是：研究者也需要 representation learning。反复阅读样本、预测标签、和领域专家核对错误，会在人的头脑中形成关于噪声、歧义、捷径与长尾的表示。Jason 用肺癌分类经历说明，先学会任务本身让他获得了后来能够提出研究问题的直觉；这并不要求每个研究者成为完整领域专家，却要求至少理解错误意味着什么。

\begin{lstlisting}[caption={Manual data inspection 的最小研究循环}]
for example in sample(training_corpus):
    guess = human_predict(example.context)
    compare(guess, example.next_token)
    label_required_signal(example)
    record_noise_shortcut_ambiguity(example)
summarize_recurring_patterns()
design_one_testable_experiment()
\end{lstlisting}

\subsection{Hyung Won 的问题：跟不上最新进展时，研究变化本身}

上一小节从数据建立直觉；现在把时间尺度拉长。Hyung Won 的标题强调 ``from the history of Transformer''，但目的不是做编年史，而是从历史中识别反复出现的驱动力。读下面几页时，应该把 ``history'' 理解为一组自然实验：算力、数据、目标和结构发生变化后，哪些方法被保留，哪些被淘汰，为什么。

\lecturefigure{hyung-slide-001.jpg}{Hyung Won Chung 的讲座主题：从 Transformer 历史塑造 AI 的未来}{Hyung Won Chung official deck, p.~1}

\lecturefigure{hyung-slide-002.jpg}{不要只追最新结果，要研究变化本身}{Hyung Won Chung official deck, p.~2}

第二页把研究焦虑重新定义成方法问题。每周追逐新模型会得到大量事实，却未必得到方向；研究 ``change itself'' 则要求把旧系统与新系统对齐，找出性能变化背后的共同变量。这样做不是为了精确预测某个产品名称，而是为了提高对研究方向的判断概率。

\teachervoice{Hyung Won 在 00:32:29--00:33:22 强调：当变化速度超过人的追踪能力时，更应该回看旧东西，画出“我们怎样走到现在”的路径。课堂重点不是记住 2024 年的模型榜单，而是练习从路径中抽象方向。}

\lecturefigure{hyung-slide-003.jpg}{研究变化的三步法：识别、理解并外推主导力量}{Hyung Won Chung official deck, p.~3}

\begin{importantbox}{Dominant driving force 的工作定义}
主导驱动力不是系统中唯一的力量，而是在当前问题尺度上解释大部分方向性的力量。研究者可以暂时忽略较小因素，以换取可分析性；但必须公开说明这是近似，并在预测失败时重新检查被忽略的因素。
\end{importantbox}

\lecturefigure{hyung-slide-004.jpg}{落笔实验：用重力说明如何忽略次要力量}{Hyung Won Chung official deck, p.~4}

对一支从静止释放的笔，若把向下方向记为正，近地面且忽略空气阻力时有
\[
y(t)=y_0+v_0t+\frac{1}{2}gt^2.
\]
其中 $y(t)$ 表示时刻 $t$ 的位置，$y_0$ 是初始位置，$v_0$ 是初速度，$g$ 是重力加速度。公式之所以可用，不是因为阻力不存在，而是因为在这个演示的精度要求下，重力足够主导且已有可靠模型。

\teachervoice{课堂在 00:35:04--00:36:21 明确提到空气阻力和复杂气动相互作用，但选择忽略它们。这个例子教的不是“现实只有一个变量”，而是“先定义精度需求，再决定哪些变量可以被近似掉”。}

\lecturefigure{hyung-slide-006.jpg}{主导力量越多、相互作用越复杂，轨迹越难预测}{Hyung Won Chung official deck, p.~6}

\begin{knowledgebox}{读图：不要把横轴误读成“变量总数”}
横轴是\textbf{主导力量的数量}，不是观测字段的数量。一个系统可以有几千个特征，但仍被少数机制支配；反过来，一个看似简单的系统若有多个强耦合反馈，也可能难以外推。这张图支持的是建模策略：先寻找方向性最强的因素，再决定是否需要加入相互作用。
\end{knowledgebox}

\subsection{共同研究循环}

前两小节分别从数据截面和历史纵轴出发，现在可以把它们合成一个循环。先看真实数据以避免脱离分布，再画曲线或历史轨迹以识别方向；随后提出最小机制解释，并主动寻找会推翻解释的反例。两位讲者的差别只在观察尺度，不在认识论结构。

\begin{table}[H]
\centering
\begin{tabularx}{\textwidth}{L{0.19\textwidth} >{\raggedright\arraybackslash}X >{\raggedright\arraybackslash}X}
\toprule
步骤 & Jason 的版本 & Hyung Won 的版本 \\
\midrule
接触证据 & 手工阅读训练样本与 prompt & 比较旧架构与当前架构 \\
寻找结构 & 分解 latent tasks 与 loss mixture & 识别 dominant driving force \\
建立曲线 & 画 compute/data/performance scaling curve & 沿时间和规模外推驱动力 \\
检查边界 & 看 metric、数据质量与缺失中间点 & 看被忽略力量、物理瓶颈与 regime change \\
形成行动 & 设计下一组规模化实验 & 添加或移除归纳偏置 \\
\bottomrule
\end{tabularx}
\caption{两场讲座共享的证据驱动研究循环。}
\end{table}

\subsection{本章小结}

本章确立了整堂课的阅读方式。Jason 要求研究者训练自己的数据直觉，Hyung Won 要求研究者训练自己的历史直觉；两者都反对从单个结果直接跳到宏大结论。接下来的技术内容将不断重复同一动作：拆解目标、读懂曲线、识别假设、说明证据不能推出什么。

